\section{Неопределённость}
\label{sec:unc}

\emph{Статус раздела.} Как и предыдущий, это замысел, а не работающий код: раздел написан до
измерений и с тех пор правился один раз. Ни одно число документа на описанной здесь
конструкции не получено.

На выходе нужно распределение, а не число, и не ради красоты: решающий слой
(раздел~\ref{sec:decision}) без него не работает вообще --- ему нужно знать не только, какую
активность модель считает наиболее вероятной, но и насколько она в этом уверена.

Неопределённость принято делить надвое, и деление это не педантизм, а рабочий инструмент.
\emph{Эпистемическая} часть --- та, что происходит от нашего незнания: похожих молекул в
обучении было мало, и модель просто не знает, что тут бывает. Она уменьшается, если добавить
данных. \emph{Алеаторная} часть --- собственный шум измерения: сколько ни добавляй молекул,
прибор точнее не станет. Различать их нужно потому, что реагировать на них надо
по-разному, и потому что трансдуктивный слой опирается ровно на первую.

Собирается распределение из трёх частей.

\emph{Ансамбль} по сидам и семействам энкодеров даёт эпистемическую часть --- ту неопределённость, которая
происходит от незнания правильной модели и падает с ростом данных. Дешёвые альтернативы, если бюджет
не позволяет держать ансамбль: SWAG или лапласовское приближение на последнем слое.

\emph{Гетероскедастичные головы} $(\sigma^-,\sigma^+)$ дают алеаторную часть --- шум прибора, который
не уменьшить никаким количеством данных.

\emph{Конформная калибровка} на разбиении, воспроизводящем дизайн теста. Обычная конформная калибровка
даёт маргинальное покрытие: в среднем по всей выборке 90 процентов заявленных интервалов накрывают истину.
Нам нужно условное: у неактивных соединений полосы широкие по построению, и если модель систематически
промахивается на активных, среднее покрытие это замаскирует. Проверяем покрытие по зонам активности
отдельно, при необходимости берём мондриановский вариант с зонами в качестве таксономии.

Разделять алеаторную и эпистемическую части нужно не для отчёта, а для трансдуктивного слоя:
эпистемическая компонента велика там, где мало похожих обучающих примеров, а тест как раз состоит
из серий аналогов вокруг известных хитов.
